\chapter{Algorithms and Programming}
\section{Summaries and Interval Scans}
\noindent\textbf{Key Ideas and Methods.}\enspace
For a data stream, first identify the summary that must be retained. Sample size, mean, and centered sum of squares suffice to update variance; merging groups also requires a correction for their different means. For interval extrema, a monotone deque retains candidates that have neither expired nor been dominated by a later element. If extending a window can only make its constraint harder to satisfy, moving its left endpoint forward can locate the longest feasible interval.

Correctness follows from the summary or deque invariant. Linear total time follows because each element enters and leaves at most once. Random simulation also requires a proof of the accepted distribution and termination; retries must use fresh independent trials.
\par\smallskip


\problem{A31}{Merging and Separating Summary Statistics}{}
Two disjoint data groups are available only through summaries $(n_A,\mu_A,M_A)$ and $(n_B,\mu_B,M_B)$, with $n_A,n_B\ge1$ and $M_A=\sum_{i\in A}(x_i-\mu_A)^2$, similarly for B. Raw observations cannot be read. Recover the combined mean, centered sum of squares, and both variance conventions. Also recover B's summary from the combined summary and A's summary. Discuss cost and numerical limitations.
\solution Put $n=n_A+n_B$ and $\delta=\mu_B-\mu_A$. Adding the two data sums gives
\[
\mu=\mu_A+\frac{n_B}{n}\delta.
\]
For group A, expand $x_i-\mu=(x_i-\mu_A)+(\mu_A-\mu)$. The cross term sums to zero because the within-group deviations sum to zero. Doing the same for B yields
\[
M=M_A+M_B+n_A(\mu_A-\mu)^2+n_B(\mu_B-\mu)^2
 =M_A+M_B+\frac{n_An_B}{n}\delta^2.
\]
The population variance is $M/n$ and the sample variance with a degrees-of-freedom correction is $M/(n-1)$. Unbiasedness of the latter additionally requires suitable iid sampling assumptions.

For separation, first find $n_B=n-n_A>0$. Then
\[
\mu_B=\frac{n\mu-n_A\mu_A}{n_B},\qquad
M_B=M-M_A-\frac{n_An_B}{n}(\mu_B-\mu_A)^2.
\]
Each operation uses constant time and a constant-size summary. In exact arithmetic, any binary merge order gives the same result: each merge recovers the actual combined moments. Floating-point rounding can make orders differ slightly. Separation can suffer severe cancellation, especially when B is small and the subtracted quantities are large. A substantially negative recovered $M_B$ signals inconsistent summaries or numerical failure and should not simply be clipped to zero.
\insight Combined variation includes both within-group variation and variation between group means. Averaging the two variances alone omits the latter.

\problem{A32}{The Longest Interval with Bounded Range}{}
Given real values $x_0,\ldots,x_{T-1}$ and $D\ge0$, find the longest nonempty interval satisfying
\[
\max_{i\le t\le j}x_t-\min_{i\le t\le j}x_t\le D.
\]
Break ties by the earliest start. Return an empty interval for empty input. Give an $O(T)$ algorithm and prove that advancing the left endpoint cannot discard an optimum.
\solution Maintain a left endpoint $\ell$, a deque of indices with decreasing values for the maximum, and another with increasing values for the minimum. On adding right endpoint $r$, remove values at most $x_r$ from the maximum deque's tail and values at least $x_r$ from the minimum deque's tail, then append $r$ to both.

While the two front values differ by more than $D$, advance $\ell$ and discard front indices smaller than the new $\ell$. The front values then give the true window extrema. Record $r-\ell+1$ only when it strictly exceeds the best length; increasing right endpoints thereby preserve the earliest start among ties.

If $[\ell,r]$ is invalid, every interval containing it is invalid. Discarding that left endpoint therefore loses no feasible interval ending now or later. When shrinking stops, $\ell$ is the earliest still-possible feasible start for $r$, so the current window is the longest feasible interval ending there. Considering every $r$ proves global optimality.

A tail index removed by a later, more extreme value is dominated: it expires no later and can never be needed as an extremum. Each index enters and leaves each deque at most once, and $\ell$ advances at most $T$ times. Time is $O(T)$ and worst-case space is $O(T)$. Every singleton is feasible when the input is nonempty.

\problem{P09}{Simulating a Die with an Unknown Biased Coin}{}
You may only toss one coin independently, with an unknown head probability $p\in(0,1)$, and have no other randomness. Design an algorithm that does not know $p$ and outputs each of $1,\ldots,6$ with exactly equal probability. Prove uniformity and almost-sure termination, and find the expected number of coin tosses used by your algorithm.
\solution First extract fair bits. Toss twice: output 0 for HT, output 1 for TH, and discard HH or TT and repeat. HT and TH each have probability $p(1-p)$, so an accepted output is fair. Successive outputs use disjoint segments of independent trials; after one segment ends, the next has the same fresh distribution. The extracted bits are consequently independent.

Let $b$ be the mean toss count per fair bit. A pair is accepted with probability $2p(1-p)>0$, giving a geometric tail, and
\[
b=2+\bigl(1-2p(1-p)\bigr)b
  =\frac1{p(1-p)}.
\]
Read three fair bits as $Z\in\{0,\ldots,7\}$. If $Z\le5$, output $Z+1$; otherwise discard the block and repeat. All eight values are equally likely, so the accepted six remain uniform. A block is accepted with probability $3/4$.

Writing $E$ for the total expected toss count, one block costs $3b$ in expectation and a rejection starts fresh. Hence
\[
E=3b+\frac14E,\qquad
\boxed{E=\frac4{p(1-p)}}.
\]
Both acceptance probabilities are positive, proving almost-sure termination. This is the cost of the stated algorithm, not a claim of optimality. Fair-bit extraction becomes expensive as $p$ approaches either endpoint.


\par\Needspace{8\baselineskip}
\subsection*{Section Exercises}
\exercise{OA31}{Rolling Mean and Variance}
Real observations $x_1,x_2,\ldots$ arrive sequentially. For a fixed window size $N\ge1$, update the mean and variance of the most recent $N$ observations whenever a new value arrives, using all observations if fewer than $N$ have arrived. Define the variance as the average squared deviation within the current window. Give an efficient algorithm, its update formulas, and its time and space complexity.

\exercise{OA32}{Sliding-Window Maxima}
Given a sequential stream of real values $x_1,x_2,\ldots$ and a positive integer $N$, output the maximum of the most recent $N$ values at each arrival, using all available values before the window fills. Design an efficient algorithm, prove its correctness, and distinguish its worst-case time per update from its amortized time.


\section{Drawdown Intervals}
\noindent\textbf{Key Ideas and Methods.}\enspace
For positive prices, fix the endpoint $j$ first. The drawdown $1-p_j/p_i$ increases with the starting price $p_i$, so maximum drawdown reduces to finding the highest price among eligible starting points. A duration limit determines the candidate window. The positivity assumption is essential to this comparison. Finding the shortest interval that reaches a threshold is a different optimization problem: later starting points with higher prices dominate earlier, lower ones.

Before discarding a candidate, justify why it cannot improve a future answer: it has expired, is dominated, or has already attained its shortest feasible interval. Arrange insertion and queries so that only indices $i<j$ are eligible. Rules for equal prices and equal objective values determine which tied interval is returned.
\par\smallskip


\problem{A33}{Maximum Drawdown with a Duration Limit}{}
Given strictly positive prices $p_0,\ldots,p_{T-1}$ and an integer $L\ge1$, allow only drawdown intervals of duration between $1$ and $L$:
\[
D=\max\left(0,\max_{1\le j-i\le L}\left(1-\frac{p_j}{p_i}\right)\right).
\]
For $T<2$, define $D=0$. If $D>0$, return an optimum, breaking ties first by the earliest end and then by the earliest start. Otherwise return zero and no interval. Design a linear-time algorithm.
\solution For a fixed end $j$, the drawdown increases with $p_i$ because $p_j>0$. Thus find the largest price in the eligible window
\[
\max(0,j-L)\le i\le j-1.
\]
Maintain a decreasing deque of candidates. Before evaluating $j$, insert $j-1$, removing only strictly smaller prices from the tail; keep equal prices to retain the earlier start. Then expire front indices smaller than $j-L$.

With front index $i$, evaluate $1-p_j/p_i$. Scan $j=1,\ldots,T-1$ and update only on a strict improvement over the best positive drawdown. This retains the earliest end; the deque's handling of equal peaks retains the earliest start. For $T<2$, return zero and no interval.

A removed tail price is dominated by a later, larger price, which is at least as good in every future eligible window and expires no earlier. The front is therefore always an eligible maximum. Optimizing separately for each end and then over all ends proves correctness. Each index enters and leaves the deque at most once, giving $O(T)$ time and $O(\min(T,L))$ space.
\pitfall The historical maximum can be too old to be eligible. A cumulative maximum alone can return an infeasible interval.

\problem{A34}{The Shortest Interval Reaching a Drawdown Threshold}{}
For strictly positive prices and $d\in(0,1)$, find the shortest interval $(i,j)$ with $i<j$ and $p_j\le(1-d)p_i$. Report failure if none exists; break length ties by the earliest end. Give an $O(T)$ algorithm and explain why maximizing the drawdown ratio does not solve this problem.
\solution Keep candidate starts in a deque with strictly decreasing prices. At end $j$, repeatedly inspect the front $i$. If
\[
p_j\le(1-d)p_i,
\]
update the shortest length with $j-i$ and remove $i$; otherwise stop this loop. Next remove tail indices whose prices are at most $p_j$, and append $j$. Querying before insertion enforces $i<j$.

If the front fails the threshold, every later candidate has a smaller price and also fails. Removing a successful front is safe: the current end is the earliest unprocessed end, and every future end would give that same start a longer interval. Tail removal is safe for a different reason. A later start with a price at least as high meets any future threshold at least as easily and produces a shorter duration.

Thus front removal records the best available duration for a start, whereas tail removal uses domination by a later, higher price. Each index enters and leaves at most once, giving $O(T)$ time and $O(T)$ space. Scan ends in order and update only on a strictly shorter length to retain the earliest end. Maximum drawdown optimizes the price ratio and may take a long time; it does not optimize duration.

\par\Needspace{12\baselineskip}
\subsection*{Section Exercises}
\exercise{OA33}{Maximum Drawdown}
Let $p_1,\ldots,p_n$ be positive prices, with $n\ge1$. Define the maximum drawdown as the largest relative decline from an earlier to a later price:
\[
 D=\max\left(0,\max_{1\le i<j\le n}\frac{p_i-p_j}{p_i}\right),
\]
with $D=0$ when $n=1$. Give an $O(n)$ algorithm, prove it correct, and state its space complexity. If $D>0$, explain how to return an interval attaining it.


\section{Algorithm Problems}
\noindent\textbf{Key Ideas and Methods.}\enspace
In dynamic programming, define what information a state retains before writing transitions and boundary values. Limits on deletions, inventory, or selected intervals must be represented in the state. Unreachable states must be distinguished from feasible states with zero reward. In counting problems, different transitions must describe disjoint sets of solutions to avoid double counting.

The choice of algorithm rests on a proved structure: binary search needs a monotone exclusion rule, greedy methods need an exchange or dominance argument, and divide-and-conquer needs a way to account for contributions across subproblems. Count all states, transitions, and data-structure operations when bounding complexity. Space compression also requires an update order that does not accidentally reuse current-round values. For randomized algorithms, distinguish distributional correctness, expected cost, and worst-case cost.
\par\smallskip


\algproblem{ALG01}{Maximum Subarray Sum with One Deletion}
Given a nonempty real array $a_0,\ldots,a_{T-1}$, choose a contiguous interval and delete at most one of its elements, retaining at least one. Find the maximum retained sum in $O(T)$ time and $O(1)$ extra space.
\solution Let $f_i$ be the best sum of an interval ending at $i$ with no deletion, and $g_i$ the best nonempty retained sum for such an interval with at most one deletion. In $g_i$, position $i$ itself may be deleted. Initialize $f_0=g_0=a_0$. For $i\ge1$,
\[
f_i=\max(a_i,f_{i-1}+a_i),\qquad
g_i=\max(a_i,g_{i-1}+a_i,f_{i-1}).
\]
The three choices in $g_i$ start afresh, retain $a_i$ and extend a previous interval, or delete $a_i$. If a solution deletes its first element, moving the interval's start right gives the same retained sequence, so no additional case is needed. All remaining solutions are covered by whether their last element is deleted.

Return $\max_i g_i$. Preserve the old $f$ when updating; the third candidate must not use its new value. Initialization uses actual elements rather than zero, preventing an invalid empty answer on an all-negative array. Induction on the two states proves optimality. Time is $O(T)$ and extra space is $O(1)$.

\algproblem{ALG02}{The Minimum Number of Jumps}
A nonempty array $a_0,\ldots,a_{T-1}$ contains nonnegative integers. From index 0, at index $i$ you may jump right by any of $1,\ldots,a_i$ positions. Find the minimum jumps to $T-1$, or report that it is unreachable. Use linear time and constant extra space; reachability is not guaranteed.
\solution Return 0 if $T=1$. Maintain the right boundary $R$ of the current reachable layer, the furthest next-layer position $F$ seen so far, and the number $k$ of completed expansions, initially all zero. Scan $i=0,\ldots,T-2$ and update
\[
F\leftarrow\max\{F,\min(T-1,i+a_i)\}.
\]
When $i=R$, the current layer has been fully scanned. If $F=R<T-1$, no new position is reachable, so report failure. Otherwise increment $k$ and set $R=F$. Return $k$ as soon as $R=T-1$.

Each jump can reach a whole forward interval. Thus positions reachable in at most $k$ jumps form a prefix starting at 0. After scanning this prefix, the union of all one-step extensions ends exactly at $F$, the boundary for at most $k+1$ jumps. Processing these layers in order proves that the first layer covering the target has minimum depth. Every index is scanned once, giving $O(T)$ time and $O(1)$ extra space.

\algproblem{ALG03}{One Trade with a Holding-Period Constraint}
Prices $p_0,\ldots,p_{T-1}$ are nonnegative, and integers $1\le h\le L$ are given. Make at most one purchase and sale, with holding constraint $h\le j-i\le L$ and sale fee $c\ge0$. Choosing no trade is allowed. Find the maximum net profit and an optimal date pair in $O(T)$ time.
\solution For sale day $j$, eligible purchase days are
\[
\max(0,j-L)\le i\le j-h.
\]
An empty range gives no feasible trade that day. Maintain an increasing price deque. On day $j$, insert $j-h$ if nonnegative, removing strictly larger tail prices, and then expire front indices smaller than $j-L$. Keeping equal prices permits the earlier purchase date.

The front minimizes eligible purchase price, so evaluate $p_j-p_i-c$. Initialize the answer to zero and record dates only on a strict improvement. Return no trade if the final value remains zero. If $T\le h$, no date pair is possible.

For a fixed sale day, the sale price and fee are fixed, making the cheapest eligible purchase optimal. A later, cheaper price dominates an older, dearer one and expires no earlier. Delayed insertion enforces the minimum holding period, and front expiration enforces the maximum. Each index enters and leaves at most once: time $O(T)$ and space $O(\min(T,L))$.

\algproblem{ALG04}{Daily Trading with Bounded Inventory}
Given nonnegative prices for $T$ days and integer inventory capacity $K\ge0$, each day buy one unit, sell one unit, or do nothing. Short selling is forbidden; each purchased or sold unit incurs fee $c\ge0$. Start with zero cash and inventory, allow negative cash, and finish with zero inventory. Maximize final cash using an $O(T(K+1))$ dynamic program.
\solution Let $V_t(q)$ be the maximum cash after the first $t$ days with $q$ units, $0\le q\le K$. Set $V_0(0)=0$ and all other initial states to $-\infty$. On the day with price $p_t$,
\[
V_{t+1}(q)=\max\left\{
\begin{array}{l}
V_t(q),\\
V_t(q-1)-p_t-c\quad(q\ge1),\\
V_t(q+1)+p_t-c\quad(q<K).
\end{array}\right\}
\]
These are no action, purchase, and sale. Every right-hand side must use the previous day's array. Updating repeatedly in place can incorrectly allow multiple units to be traded on one day.

Every legal strategy belongs to exactly one last-action case. Removing that action gives the corresponding prior state, and appending it to any finite prior state is feasible. Induction proves both the upper bound and attainability. Return $V_T(0)\ge0$, since no action is always feasible. Time is $O(T(K+1))$ and rolling-array space is $O(K+1)$; stored choices recover a schedule. If $K=0$ or $T=0$, return zero. Allowing negative cash is a material assumption; a cash constraint would require further feasibility checks.

\algproblem{ALG05}{Counting Paths through at Most One Obstacle}
An $R\times C$ binary grid, with $R,C\ge1$, marks obstacles by 1. Move only right or down from the upper-left to the lower-right cell, visiting at most one obstacle, with endpoints counted by the same rule. Find the number of paths modulo an integer $M\ge2$ in $O(RC)$ time and $O(C)$ space.
\solution Let $A_{i,j}$ and $B_{i,j}$ count paths reaching $(i,j)$ after visiting zero and exactly one obstacle. Sum the corresponding counts from above and from the left, obtaining $a,b$. At a zero cell, set $(A,B)=(a,b)$; at an obstacle, set $(A,B)=(0,a)$. Reduce every addition modulo $M$.

Initialize the start separately as $(1,0)$ if it is zero and $(0,1)$ otherwise. Other out-of-grid predecessors contribute zero. Return $A_{R-1,C-1}+B_{R-1,C-1}\pmod M$.

Every path's penultimate cell is either above or left, disjoint cases. Its obstacle count then changes exactly according to the current cell. Thus the recurrence neither misses nor double-counts a path. Scan each row left to right with two length-$C$ arrays: the old current column stores the upper predecessor, and the updated previous column stores the left predecessor. Read both before overwriting. For $R,C\ge1$, the same initialization covers a one-cell grid.

\algproblem{ALG06}{A Shortest Common Supersequence}
Given strings $a,b$, possibly empty, construct a shortest string $s$ containing each as a subsequence. A character of $s$ may serve one character from each input simultaneously. Find the minimum length, recover a solution, and prove correctness.
\solution Let $D_{i,j}$ be the shortest common supersequence length for prefixes of lengths $i,j$. Initialize $D_{i,0}=i$ and $D_{0,j}=j$. Index characters from 1. If the last characters agree, share one final character:
\[
D_{i,j}=D_{i-1,j-1}+1\quad(a_i=b_j).
\]
If they differ, the final character removes the last character of only one input, giving
\[
D_{i,j}=1+\min(D_{i-1,j},D_{i,j-1}).
\]
When the last characters agree, their two matches in any supersequence can both be moved to the later matched position, so sharing that final position is harmless. For different last characters, the two alternatives exhaust an optimum. Induction on prefix lengths proves the recurrence.

Backtrack from $(|a|,|b|)$. On equal characters record that character and move diagonally; otherwise follow a minimum predecessor and record the removed last character. At a boundary, append the remaining prefix in reverse order, then reverse the entire record. Time and storage are $O((|a|+1)(|b|+1))$.

If $L$ is the longest common subsequence length, the answer is also $|a|+|b|-L$. Shared matched positions form a common subsequence, so at most $L$ positions can be shared. Conversely, align a longest common subsequence and insert the intervening characters of both strings to attain this bound.

\algproblem{ALG07}{Counting and Sampling Interleaving Assignments}
Given strings $a,b,t$ with $|t|=|a|+|b|$, assign every position of $t$ to one input while preserving the order within each input. Distinct source assignments count separately even when their characters coincide. Count legal assignments and, when the count is positive, sample one uniformly.
\solution Let $C_{i,j}$ count completions after using the first $i$ characters of $a$ and $j$ of $b$. Put $m=|a|$, $n=|b|$, and $C_{m,n}=1$. With characters indexed from 1, compute backwards:
\[
C_{i,j}=\ind_{\{i<m,\ a_{i+1}=t_{i+j+1}\}}C_{i+1,j}
       +\ind_{\{j<n,\ b_{j+1}=t_{i+j+1}\}}C_{i,j+1}.
\]
Out-of-range branches contribute zero. The two branches have different sources, so their counts add even if the characters agree. The total is $C_{0,0}$; a length mismatch gives zero immediately.

At a state with positive count, select the next source with probability equal to that branch's count divided by $C_{i,j}$, then move to its successor. For any complete assignment, multiplying these probabilities telescopes to $1/C_{0,0}$, proving uniformity. The boundary also covers the unique empty assignment.

There are $O((m+1)(n+1))$ states and recovery takes $m+n$ steps. Counts are exact integers; the state-operation bound does not assume constant bit cost for arbitrarily large integer additions. Counts reduced modulo an integer cannot supply exact sampling weights.

\algproblem{ALG08}{Weighted Interval Selection with a Cardinality Limit}
There are $n$ tasks. Task $i$ occupies a half-open interval $[s_i,e_i)$ with $s_i<e_i$ and has possibly negative reward $w_i$. Given an integer $K\ge0$, choose at most $K$ pairwise nonoverlapping tasks; one may start exactly when another ends. Find the maximum reward in $O(n\log n+nK)$ time.
\solution Sort by end time and let $p(i)$ be the largest predecessor index with $e_j\le s_i$, or 0 if none exists. Since $s_i<e_i$, every such predecessor appears before $i$. Binary search on sorted end times finds all predecessors in $O(n\log n)$.

Let $F(i,k)$ be the best reward from the first $i$ tasks using at most $k$. Initialize $F(0,k)=F(i,0)=0$ and use
\[
F(i,k)=\max\{F(i-1,k),\ w_i+F(p(i),k-1)\}.
\]
Excluding $i$ gives the first case. If $i$ is included, all other chosen tasks must end before it starts, giving the second. Conversely, an optimal compatible predecessor set can be joined to $i$, proving attainability as well as the bound.

Replace $K$ by $\min(K,n)$ and return $F(n,K)$. Negative-reward tasks can be skipped, so zero is always attainable. A full $O(n(K+1))$ table supports backtracking. If only the value is needed, process $k$ by layers with two arrays of length $n+1$.

\algproblem{ALG09}{Scheduling Tasks with Individual Cooldowns}
Day $i$ offers a task available only that day, with real reward $r_i$. After performing it, rest for the next $L_i\ge0$ complete days, where $L_i$ is an integer; no other task may be performed during rest. Tasks may be skipped, for $i=0,\ldots,T-1$. Maximize total reward and recover the selected dates in $O(T)$ time.
\solution Let $V_i$ be the best reward starting on day $i$, with $V_T=0$. After choosing today's task, the next eligible day is
\[
q_i=\min(T,i+L_i+1).
\]
Compute backwards:
\[
V_i=\max\{V_{i+1},r_i+V_{q_i}\},\qquad i=T-1,\ldots,0.
\]
The alternatives skip today or perform the task and resume after its rest period. They cover every legal schedule, and joining either action to an optimal successor schedule is feasible. Backward induction proves optimality.

Store the chosen branch. Starting at 0, advance one day after a skip; after a chosen task, output $i$ and jump to $q_i$. Indices strictly increase, so recovery visits at most $T$ states. Time and space are $O(T)$. Skipping ensures that nonpositive rewards need not reduce the answer; empty input gives zero. Variable successor gaps prevent a fixed-size recent-state buffer in general.

\algproblem{ALG10}{Repairing Parentheses with the Fewest Insertions}
Given a string containing only opening and closing parentheses, insert characters without deleting or changing any original character. Minimize the insertions needed to make the entire string valid, and output a repair attaining the minimum. Use linear time.
\solution Scan left to right with the number $b$ of unmatched opening parentheses. Output an opening parenthesis unchanged and increment $b$. On a closing parenthesis, if $b>0$, output it and decrement $b$; otherwise insert an opening parenthesis immediately before it and count one insertion. At the end append $b$ closing parentheses.

Every output prefix has nonnegative balance and the final balance is zero, so the repair is valid and retains all original characters in order.

For optimality, let $m\le0$ be the minimum original prefix balance (including the empty prefix) and $s$ the final original balance. Every repair needs at least $L=-m$ inserted openings, or the most deficient prefix remains invalid. If $\ell,r$ openings and closings are inserted, final balance requires $r=s+\ell$. The total is therefore $s+2\ell\ge s+2L$. The algorithm inserts exactly $L$ openings, leaving balance $s+L$, and appends that many closings. It attains the lower bound in $O(T)$ time and constant space apart from its output.

\algproblem{ALG11}{Water Retention on a Two-Dimensional Terrain}
A rectangular grid has nonnegative ground height $h_v$ and unit area at each cell. Neighbors share an edge. Water can drain freely outside, and every boundary cell connects to the exterior. After sustained rain reaches equilibrium, find the total retained water. Give an algorithm and explain why row and column maxima alone are insufficient.
\solution Define the escape level of cell $v$ by
\[
\ell_v=\min_{\gamma:\,v\text{ to boundary}}\ \max_{u\in\gamma}h_u.
\]
Paths include their endpoints, so $\ell_v\ge h_v$. The largest height on a path is its drainage threshold; taking the lowest threshold among all escape paths gives the stable water level. Total water is $\sum_v(\ell_v-h_v)$.

Solve this minimax path problem with a min-heap. Initialize every boundary cell at distance $h_v$ and interior cells at $+\infty$. Extract the unsettled cell $v$ with smallest distance and relax each neighbor $u$ by
\[
d_u\leftarrow\min\{d_u,\max(d_v,h_u)\}.
\]
Ignore obsolete heap entries. Extending a path never reduces its maximum height, so the usual shortest-path greedy argument applies. More explicitly, a lower escape path would contain a first unsettled cell adjacent to a settled one; relaxation from that settled cell would already have supplied a no-larger candidate, contradicting the minimum heap extraction.

There are only a constant number of edges per cell, giving $O(RC\log(RC+1))$ time and $O(RC)$ space. With one row or one column, every cell is on the boundary and retention is zero. Drainage paths may turn, so barriers measured only along straight rows or columns can overestimate retention.

\algproblem{ALG12}{Largest and Total All-One Squares}
For a binary matrix, find both the largest area of an axis-aligned all-one square and the total number of such squares. Different locations or side lengths count separately. Use $O(RC)$ time and $O(C)$ extra space.
\solution Let $d_{i,j}$ be the largest side length of an all-one square with lower-right corner $(i,j)$. A zero cell gives $d_{i,j}=0$; a one cell gives
\[
d_{i,j}=1+\min(d_{i-1,j},d_{i,j-1},d_{i-1,j-1}),
\]
with out-of-grid values zero. A side-$k$ square requires side-$(k-1)$ squares at all three neighboring corners, proving necessity. Conversely, those three smaller squares together with the current cell cover the entire side-$k$ square, proving sufficiency.

The largest area is $(\max d_{i,j})^2$. At a fixed lower-right corner, feasible side lengths are exactly $1,\ldots,d_{i,j}$, so the total count is
\[
\sum_{i,j}d_{i,j}.
\]
Every square has one lower-right corner, avoiding double counting. Scan rows with a length-$C$ array and one saved old upper-left value: the old current column is the upper neighbor and the updated previous column is the left neighbor. Both answers are zero for an all-zero or empty matrix.

\problem{CA01}{Weighted Sampling from a Stream}{}
Items arrive sequentially with positive weights $w_i$. The total number of items is unknown. Only one candidate item and a constant number of numerical values may be stored. Design an algorithm such that after any $t\ge1$ arrivals, item $i$ is the candidate with probability $w_i/\sum_{j=1}^t w_j$. Each update may use an independent continuous uniform random variable. Prove the claim and explain why replacing the candidate whenever the new item is heavier fails.
\solution Maintain the cumulative weight $W$. Select the first item and set $W=w_1$. When item $t$ arrives, first update $W\leftarrow W+w_t$, then replace the candidate with probability $w_t/W$. Otherwise retain it. Each update takes constant arithmetic time and stores only the candidate identifier and cumulative weight.

For the inductive proof, let $W_{t-1}$ denote the previous total. An old item $i<t$ survives precisely when it was selected previously and the new update does not replace it. Independence of the fresh random variable gives
\[
\Prob(\text{candidate}=i)
=\frac{w_i}{W_{t-1}}\left(1-\frac{w_t}{W_t}\right)
=\frac{w_i}{W_t}.
\]
The new item is selected with probability $w_t/W_t$ directly. The base case $t=1$ is immediate, proving the invariant at every fixed length. It also applies when the ending length is fixed externally or independent of the sampling randomness.

For weights $1,2,3$, the second update selects item 2 with probability $2/3$ and the third replaces with probability $1/2$. Final probabilities are $1/6,1/3,1/2$. Deterministic replacement by a heavier item always retains the last item in a strictly increasing stream, so it cannot have the required distribution.

The proof applies to a fixed sequence of weights. Random weights independent of the sampling randomness can be handled by conditioning on the entire sequence. If future weights adapt to publicly revealed candidate choices, that conditioning argument is no longer valid without further analysis. The space bound counts stored values; their exact bit lengths may grow.
\insight The candidate represents all past weight through a probability invariant. A new item competes against the entire accumulated weight, not merely against the weight of the current representative.

\problem{CA02}{The Median of Two Sorted Arrays}{}
Two nondecreasing real arrays $A,B$ have lengths $m,n$ with $m+n>0$. Random access is allowed, but fully merging the arrays is not. Find their combined median, averaging the central two elements when the total length is even. Give an $O(\log(\min\{m,n\}+2))$-time, $O(1)$-extra-space algorithm covering empty arrays, ties, and very unequal lengths.
\solution Swap the arrays if necessary so that $m\le n$. If $m=0$, read the central one or two elements of $B$ directly. Otherwise put
\[
h=\left\lfloor\frac{m+n+1}{2}\right\rfloor.
\]
Take $i$ elements from the beginning of $A$ and $j=h-i$ from the beginning of $B$ as the left half. The feasible search range is
\[
\max(0,h-n)\le i\le\min(m,h).
\]
The four boundary elements are
\[
\begin{array}{ll}
L_A=A_{i-1}\ (i>0),&R_A=A_i\ (i<m),\\
L_B=B_{j-1}\ (j>0),&R_B=B_j\ (j<n).
\end{array}
\]
Use $-\infty$ for a missing left boundary and $+\infty$ for a missing right boundary; these sentinels are used only in comparisons.

A partition is correct exactly when $L_A\le R_B$ and $L_B\le R_A$. Internal sortedness then guarantees that every left-half element is at most every right-half element. If $L_A>R_B$, too many elements came from $A$, so decrease $i$. Larger $i$ cannot work, since $L_A$ is nondecreasing and $R_B$ nonincreasing as $i$ grows. If $L_B>R_A$, increase $i$ by the symmetric argument. Binary search therefore applies.

A correct partition exists: conceptually sort all elements, using any fixed stable order for ties. The first $h$ elements form a prefix of each original array. Every elimination preserves at least one correct partition, so the search terminates. For odd $m+n$, return $\max(L_A,L_B)$. For even $m+n$, return
\[
\frac{\max(L_A,L_B)+\min(R_A,R_B)}2.
\]
In the even, nonempty case both halves contain elements, so the final average uses finite data values. Ties may permit several partitions but all yield the same median. There are at most $m+1$ candidate integers, proving the time and space bounds.
\insight Search for the number of elements contributed by the shorter array to the left half. Monotonicity of the boundary tests replaces the need to merge the data.

\problem{CA03}{Inversions and Minimum Adjacent Swaps}{}
A real array $a_1,\ldots,a_n$ may contain ties. An inversion is a pair $i<j$ with $a_i>a_j$. Count inversions in $O(n\log n)$ time and prove that their number equals the minimum number of adjacent swaps needed to make the array nondecreasing. Evaluate $(3,1,2,1)$ and explain how ties are handled.
\solution Augment merge sort with an inversion count. Recursively sort each half and count its internal inversions. During merging, output the left element whenever it is at most the right element. Otherwise output the right element and add the number of unconsumed left elements: all of them are strictly larger. Each cross-half inversion is counted exactly when its right element is output first, so none is missed or duplicated.

Each merge level costs $O(n)$ and there are $O(\log n)$ levels. A merge buffer uses $O(n)$ additional space. Empty and one-element arrays have count zero. Taking the left element first on equality avoids counting ties as inversions. The counter must accommodate values up to $n(n-1)/2$.

For optimality of the swap count, an adjacent swap leaves the total contribution of every other element unchanged: that element lies on the same side of both swapped positions. Only the swapped pair can change its inversion status, so one swap decreases the total count by at most one. Any sorting sequence therefore needs at least the initial inversion count in swaps.

Conversely, any unsorted array contains a strictly decreasing adjacent pair; otherwise all adjacent comparisons would already imply sortedness. Swapping such a pair decreases the inversion count by exactly one. Repeating sorts the array in exactly the initial inversion count, attaining the lower bound. For the example, the inverted index pairs are
\[
(1,2),(1,3),(1,4),(3,4),
\]
so both the inversion count and the minimum adjacent-swap count are $4$.
\insight Divide and conquer assigns each pair a unique counting step. The same count is a potential function that certifies the optimal number of swaps.


\pagebreak
\section{Exercises}
\noindent\textbf{Key Ideas and Methods.}\enspace
These exercises test what can be recovered from a summary, when a window is feasible, and why candidates may be discarded. For variance, distinguish the population denominator from the degrees-of-freedom correction. For interval problems, specify indices, duration, and tie rules. A feasible interval establishes only feasibility; optimality also requires ruling out a better one. A linear-time argument must count how often each candidate is processed over the entire scan.
\par\smallskip


\exercise{AG4-1}{Recovering a Merged Group}
Two disjoint groups have a combined count of 5, mean 4, and centered sum of squares 26. Group A is $(1,4,7)$. Without reading group B, recover its count, mean, and centered sum of squares, and both combined variance conventions. Can B's two observations be recovered uniquely?

\exercise{AG4-2}{Endpoints of a Longest Stable Interval}
For $(4,2,3,5,8,7,6)$, find the longest interval whose range is at most 3. Use zero-based indices and the earliest start among ties. Give the result and a short independent argument excluding longer intervals.

\exercise{AG4-3}{Which Peak Does the Duration Limit Exclude?}
Prices are $(120,80,100,70,90)$. Find maximum drawdown with duration at most 2 and with no duration limit, giving zero-based endpoints. Explain why the historical maximum cannot be used directly for the first problem.

\exercise{AG4-4}{The Shortest Time to a Quarter Drawdown}
Prices are $(100,94,80,79,60)$ and the threshold is $d=1/4$. Find the shortest qualifying drawdown interval. Exclude shorter possibilities and explain why a successful starting index can be discarded permanently during an end scan.
